\chapter{Study Scenario}
\label{app:scenario}

The scenario reproduced below was presented to all participants in German, regardless of interaction mode (Form, Ask, Chat). The English translation is provided for the international reader; the German original is the version actually shown.

\section{Original (German)}
\label{app:scenario-de}

\begin{verbatim}
Szenario: Mitarbeiterverwaltung (Arbeitszeit/Urlaub) & Parkplatzberechtigung

Sie sind Mitarbeitender namens A.I. und arbeiten für die APOR GmbH, ein
Softwareunternehmen mit ca. 50 Mitarbeitenden in Ulm. Sie handeln im Namen
des Unternehmens und sind im Szenario sowohl Verantwortlicher (Data
Controller) als auch Datenschutzbeauftragte:r (DPO/DSB).

Zur Organisation:

Rolle: Data Controller, DPO
Adresse: Mainstraße 67, Ulm
E-Mail: mail@apor.eu
Telefon: 012345678
Vertreter: (Sie), A. I., ai@apor.eu
DPO: Das sind auch Sie

Verarbeitung
Das Unternehmen verarbeitet personenbezogene Daten, um Mitarbeitende vom
Eintritt bis zum Austritt zu verwalten. Im Fokus stehen:

- Anlegen neuer Mitarbeitender
- Dokumentation von Arbeitszeiten
- Verwaltung von Urlaub und Abwesenheiten

Zusätzlich verwaltet APOR für berechtigte Mitarbeitende eine
Parkplatzberechtigung für einen Firmenparkplatz mit Schranke. Damit
berechtigte Mitarbeitende den Parkplatz nutzen können, wird ein geeignetes
Merkmal zur Zuordnung der Berechtigung hinterlegt (z. B. ein
Fahrzeugkennzeichen).

APOR arbeitet mit intern gehostetem Microsoft Office sowie internen Servern
(keine externe Cloud). Krankentage oder medizinische Angaben sind für diese
Verarbeitung nicht erforderlich.

Für die Verwaltung sind technische und organisatorische Maßnahmen
vorgesehen, wie z.B. Zugriffsbeschränkung nach Rollen (Need-to-know),
Passwortschutz (ggf. MFA), regelmäßige Backups, Zutrittskontrolle zu
Serverraum/Archiv und Vertraulichkeitsverpflichtung der berechtigten
Mitarbeitenden.

Alle in diesem Zusammenhang gespeicherten Informationen werden spätestens
3 Monate nach Austritt der jeweiligen Person gelöscht.
\end{verbatim}

\section{English Translation}
\label{app:scenario-en}

\begin{verbatim}
Scenario: Employee Administration (Working Time / Leave) & Parking
Authorisation

You are an employee named A.I., working for APOR GmbH, a software company
with approximately 50 employees in Ulm. You act on behalf of the company
and, in this scenario, are both the Data Controller and the Data
Protection Officer (DPO).

Organisation details:

Role: Data Controller, DPO
Address: Mainstrasse 67, Ulm
E-mail: mail@apor.eu
Phone: 012345678
Representative: (You), A. I., ai@apor.eu
DPO: Also you

Processing
The company processes personal data in order to administer employees from
onboarding to exit. The focus is on:

- Onboarding new employees
- Documentation of working hours
- Administration of leave and absences

In addition, APOR manages a parking authorisation for entitled employees
for a company car park with a barrier. To allow entitled employees to use
the car park, a suitable feature for assigning the authorisation is
recorded (e.g. a vehicle licence plate).

APOR uses internally hosted Microsoft Office and internal servers (no
external cloud). Sick days or medical information are not required for
this processing.

For administration, technical and organisational measures (TOMs) are
provided, such as role-based access restriction (need-to-know), password
protection (where appropriate, MFA), regular backups, physical access
control to the server room / archive, and a confidentiality obligation for
authorised employees.

All information stored in this context is deleted no later than 3 months
after the person's departure.
\end{verbatim}


\chapter{LLM-as-a-Judge Configuration}
\label{app:judge}

This appendix documents the inputs to the LLM-as-a-judge evaluation layer that scored the 113 retained ROPA documents. It contains the verbatim system and user prompts, the five-dimension rubric, and the model configuration needed to reproduce the evaluation.

\section{System and User Prompt}
\label{app:judge-prompt}

The system prompt below is sent verbatim as the \texttt{system} role on every evaluation call. It is taken from \texttt{python/judge/prompt.json} (key \texttt{judge.system}).

\begin{verbatim}
You are an expert GDPR compliance reviewer evaluating Records of
Processing Activities (ROPA) documents under Article 30 GDPR. The
documents are written by novices (non-experts, no prior GDPR training)
who described the same processing scenario using different LLM-assisted
tools. Your job is to judge each ROPA document against (a) the provided
scenario and (b) GDPR Article 30 requirements. Be strict but fair. Do not
reward verbosity. Do not penalise a document for merely restating
scenario facts. Respond with valid JSON only - no prose, no Markdown
fences, no commentary before or after the JSON object.
\end{verbatim}

The accompanying user-message template is sent on every call with \texttt{\{scenario\}} (the German original from Appendix~\ref{app:scenario-de}) and \texttt{\{document\}} (the participant ROPA document) substituted in:

\begin{verbatim}
You will evaluate one ROPA document against the scenario it is supposed
to cover.

=== SCENARIO (verbatim, German) ===
{scenario}
=== END SCENARIO ===

=== ROPA DOCUMENT TO EVALUATE ===
{document}
=== END DOCUMENT ===

Score the document on each of the following five metrics using a 1-5
Likert scale, and give a short rationale (max 2 sentences, English) for
each score. The scale for all metrics is:
1 = very poor, 2 = poor, 3 = adequate, 4 = good, 5 = excellent.

Metrics and what to assess:
- completeness: Coverage of Art. 30(1) GDPR required elements -
  controller identity and contact, DPO, purposes of processing, legal
  basis, categories of data subjects and personal data, categories of
  recipients, retention periods, and technical/organisational measures
  (TOMs). Missing elements lower the score.
- scenario_faithfulness: How accurately the document reflects the
  concrete facts of THIS scenario - APOR GmbH, Ulm (Mainstrasse 67), ~50
  employees, internally hosted Microsoft Office and internal servers (NO
  external cloud), no medical/sickness data, 3-month deletion after
  employee departure, parking lot authorisation using a distinguishing
  feature such as a licence plate. Generic ROPA boilerplate that ignores
  these specifics lowers the score.
- legal_correctness: Correct use of GDPR terminology and correct citation
  of legal basis (typically Art. 6(1)(b) performance of contract for
  employment data, Art. 6(1)(c) legal obligation where applicable, Art.
  6(1)(f) legitimate interest for parking authorisation). Wrong articles,
  invented legal bases, or confusing controller/processor/DPO roles lower
  the score.
- hallucination_inverse: Inverse score for invented facts not supported
  by the scenario (e.g. invented third-party processors, invented data
  categories like medical data, invented addresses, invented retention
  periods different from 3 months). 5 = no hallucinations, 1 = many
  invented facts.
- overall: Holistic judgement of how useful this document would be as a
  real Art. 30 ROPA record for this scenario, taking the four dimensions
  above into account.

Return a single JSON object with exactly this shape:
{
  "completeness":           {"score": <int 1-5>, "rationale": "<string>"},
  "scenario_faithfulness":  {"score": <int 1-5>, "rationale": "<string>"},
  "legal_correctness":      {"score": <int 1-5>, "rationale": "<string>"},
  "hallucination_inverse":  {"score": <int 1-5>, "rationale": "<string>"},
  "overall":                {"score": <int 1-5>, "rationale": "<string>"}
}
\end{verbatim}

\section{Rubric --- Five Dimensions}
\label{app:judge-rubric}

All five dimensions share the same 1--5 Likert anchor: $1 = $ very poor, $2 = $ poor, $3 = $ adequate, $4 = $ good, $5 = $ excellent. The fourth dimension, \texttt{hallucination\_inverse}, is inverted internally so that for all five dimensions higher equals better. For each dimension the judge returns both an integer score and a short English rationale (maximum two sentences).

\begin{table}[htbp]
\centering
\small
\begin{tabularx}{\textwidth}{lX}
\toprule
\textbf{Dimension} & \textbf{What is assessed} \\
\midrule
Completeness & Coverage of Art.~30(1) GDPR required elements: controller identity and contact, DPO, purposes of processing, legal basis, categories of data subjects and personal data, categories of recipients, retention periods, and technical/organisational measures (TOMs). Missing elements lower the score. \\
\addlinespace
Scenario Faithfulness & Accuracy with respect to the concrete facts of this scenario: APOR GmbH, Ulm (Mainstraße~67), $\sim$50 employees, internally hosted Microsoft Office and internal servers (no external cloud), no medical/sickness data, 3-month deletion after employee departure, parking authorisation using a distinguishing feature such as a licence plate. Generic ROPA boilerplate lowers the score. \\
\addlinespace
Legal Correctness & Correct use of GDPR terminology and correct citation of legal basis (typically Art.~6(1)(b) for employment data, Art.~6(1)(c) where applicable, Art.~6(1)(f) for parking authorisation). Wrong articles, invented legal bases, or confused controller/processor/DPO roles lower the score. \\
\addlinespace
Hallucination (inverse) & Inverse score for invented facts not supported by the scenario (e.g. invented third-party processors, invented categories such as medical data, invented addresses, retention periods other than 3~months). 5~=~no hallucinations, 1~=~many invented facts. \\
\addlinespace
Overall & Holistic judgement of how useful the document would be as a real Art.~30 ROPA record for this scenario, taking the four dimensions above into account. \\
\bottomrule
\end{tabularx}
\caption{The five judge dimensions. All share the 1--5 Likert anchor; higher is better for every dimension after the hallucination inversion.}
\label{tab:judge-rubric}
\end{table}

\section{Model Configuration and Limitations}
\label{app:judge-config}

The configuration in Table~\ref{tab:judge-config} is taken from the \texttt{judge} block of \texttt{python/judge/prompt.\-json} and the corresponding cells of \texttt{ROPAgen\_JUDGE.ipynb}. It is the minimal information required to reproduce the evaluation layer.

\begin{table}[htbp]
\centering
\small
\begin{tabularx}{\textwidth}{lX}
\toprule
\textbf{Setting} & \textbf{Value} \\
\midrule
Model & \texttt{google/gemini-3.1-pro-preview} (via OpenRouter) \\
Temperature & \texttt{0.0} (deterministic decoding) \\
Output format & JSON object, enforced via OpenRouter \texttt{response\_format = json\_schema} with \texttt{strict: true} \\
Schema name & \texttt{ropa\_judge\_result} \\
Required keys & \texttt{completeness}, \texttt{scenario\_faithfulness}, \texttt{legal\_correctness}, \texttt{hallucination\_inverse}, \texttt{overall} \\
Per-key shape & \texttt{\{"score": int in [1,5], "rationale": string\}} \\
Additional properties & Disallowed at every object level \\
Documents evaluated & 113 ROPA documents (Form~37, Ask~37, Chat~39) \\
Caching & Per-document JSON result cached on disk; reused on re-runs \\
Retry behaviour & Schema-validation or transient API failures retried; persistent failures logged and skipped \\
Logging & Each call logs model, latency, prompt and completion token counts, and raw response payload \\
\bottomrule
\end{tabularx}
\caption{LLM-as-a-judge model configuration.}
\label{tab:judge-config}
\end{table}

\paragraph{Known limitations.} The following limitations of the judge layer are stated for full transparency and discussed further in the Discussion chapter:

\begin{itemize}
\item \textbf{Single-judge bias.} Only one judge model (Gemini~3.1 Pro Preview) operationalises the rubric; inter-rater agreement with a human GDPR expert is not established.
\item \textbf{Model-family bias.} Documents are produced by Mistral Large~3 inside ropagen, while the judge belongs to a different model family (Google). This reduces but does not eliminate stylistic-preference bias.
\item \textbf{Single-scenario bias.} All scores are conditional on the one APOR GmbH scenario; generalisation to other ROPA contexts is not claimed.
\item \textbf{Reference-free.} Unlike the NLP metric layer, the judge does not compare against an expert reference document; it scores against the scenario and Art.~30 directly.
\end{itemize}


\chapter{Full Inferential Tables}
\label{app:stats}

This appendix consolidates the canonical sample-size derivation and the full inferential statistics for both the survey and the NLP/judge layer. The tables in the Evaluation chapter draw selectively from these; the appendix is reproduced for completeness and reproducibility.

\section{Sample-Size Derivation}
\label{app:samplesizes}

All sample sizes used in the thesis are derived from the raw data files committed to the repository: \texttt{python/survey/docs/data.csv} for the survey (one row per participant) and \texttt{python/metrics/output/all\_scores\_combined.csv} for the NLP and judge scores (one row per AI-generated document).

\begin{table}[htbp]
\centering
\small
\begin{tabular}{lr}
\toprule
\textbf{Survey subset} & \textbf{n} \\
\midrule
Recruited participants (raw) & 73 \\
Participants with at least one valid session ($p\_0001 \geq 1$) & 70 \\
Form responses (mode-level) & 69 \\
Ask responses (mode-level) & 69 \\
Chat responses (mode-level) & 69 \\
Matched-triple participants (valid response in all three modes) & 69 \\
\bottomrule
\end{tabular}
\caption{Survey sample sizes after consent and quality filtering.}
\label{tab:n-survey}
\end{table}

\begin{table}[htbp]
\centering
\small
\begin{tabular}{lr}
\toprule
\textbf{NLP / judge subset} & \textbf{n} \\
\midrule
Total documents with parseable NLP scores and judge ratings & 113 \\
Form documents & 37 \\
Ask documents & 37 \\
Chat documents & 39 \\
Unique participants represented & 47 \\
Participants with at least one document in each mode & 28 \\
Strict matched triple (exactly one document per mode) & 26 \\
\bottomrule
\end{tabular}
\caption{NLP and judge sample sizes. The strict $n=26$ set is used for all within-subjects tests (Friedman, Wilcoxon signed-rank); the unbalanced $n=113$ is used for between-groups sensitivity checks.}
\label{tab:n-nlp}
\end{table}

\begin{table}[htbp]
\centering
\small
\begin{tabular}{lr}
\toprule
\textbf{Survey $\times$ NLP join} & \textbf{n} \\
\midrule
Joined rows (document with matching survey response) & 106 \\
Form joined rows & 34 \\
Ask joined rows & 34 \\
Chat joined rows & 38 \\
Unique participants in the joined sample & 44 \\
Participants with survey $\times$ NLP coverage in all three modes & 27 \\
\bottomrule
\end{tabular}
\caption{Triangulated samples used for confidence--quality analyses. The joined $n=106$ is the canonical sample for the rank-inversion table; the within-subjects $n=27$ is used for paired confidence~$\times$~quality correlations.}
\label{tab:n-triangulation}
\end{table}

\section{Survey Inferential Tables}
\label{app:stats-survey}

All survey tests use the matched-triple $n=69$ within-subjects subset. Omnibus is the Friedman test (non-parametric repeated-measures); post-hoc is pairwise Wilcoxon signed-rank with Bonferroni-corrected $\alpha = 0.0167$. Effect-size labels: $|r| < 0.1$ negligible, $0.1$--$0.3$ small, $0.3$--$0.5$ medium, $\geq 0.5$ large. For Kendall's~$W$: $0$--$0.1$ weak, $0.1$--$0.3$ moderate, $0.3$--$0.5$ strong.

\begin{table}[htbp]
\centering
\small
\begin{tabular}{lrrrrrl}
\toprule
\textbf{Measure} & \textbf{n} & $\boldsymbol{\chi^2}$ & \textbf{df} & \textbf{p} & $\boldsymbol{W}$ & \textbf{Sig.} \\
\midrule
SUS (0--100, higher = better)        & 69 & 5.51  & 2 & 0.0636 & 0.040 & No \\
R-TLX composite (1--21)              & 69 & 5.24  & 2 & 0.0728 & 0.038 & No \\
R-TLX Mental Demand                  & 69 & 3.13  & 2 & 0.2086 & 0.023 & No \\
R-TLX Physical Demand                & 69 & 2.08  & 2 & 0.3535 & 0.015 & No \\
R-TLX Temporal Demand                & 69 & 12.77 & 2 & 0.0017 & 0.093 & \textbf{Yes} \\
R-TLX Effort                         & 69 & 3.78  & 2 & 0.1511 & 0.027 & No \\
R-TLX Frustration                    & 69 & 8.93  & 2 & 0.0115 & 0.065 & \textbf{Yes} \\
R-TLX Performance (positive-coded)   & 69 & 2.33  & 2 & 0.3114 & 0.017 & No \\
AI Confidence (item \texttt{5\_ai})  & 69 & 3.76  & 2 & 0.1524 & 0.027 & No \\
AI Support (6-item composite)        & 69 & 1.05  & 2 & 0.5919 & 0.008 & No \\
AI Reuse intention (item \texttt{6\_ai}) & 69 & 11.12 & 2 & 0.0039 & 0.081 & \textbf{Yes} \\
\bottomrule
\end{tabular}
\caption{Friedman omnibus across all survey measures ($n = 69$ matched triple).}
\label{tab:survey-friedman}
\end{table}

\begin{table}[htbp]
\centering
\small
\begin{tabular}{llrrrrl}
\toprule
\textbf{Measure} & \textbf{Pair} & \textbf{n} & \textbf{W} & \textbf{p} & \textbf{r} & \textbf{Sig.\ (Bonf.)} \\
\midrule
SUS              & Form vs Ask  & 60 & 578.5 & 0.0131  & $+0.368$ & \textbf{Yes} \\
SUS              & Form vs Chat & 58 & 549.5 & 0.0177  & $+0.358$ & No \\
SUS              & Ask vs Chat  & 61 & 885.5 & 0.6660  & $-0.063$ & No \\
R-TLX comp.\     & Form vs Ask  & 65 & 684.0 & 0.0111  & $-0.362$ & \textbf{Yes} \\
R-TLX comp.\     & Form vs Chat & 65 & 636.0 & 0.0043  & $-0.407$ & \textbf{Yes} \\
R-TLX comp.\     & Ask vs Chat  & 63 & 995.5 & 0.9318  & $+0.012$ & No \\
Temporal Demand  & Form vs Ask  & 43 & 237.0 & 0.0038  & $-0.499$ & \textbf{Yes} \\
Temporal Demand  & Form vs Chat & 46 & 233.0 & $<$0.001 & $-0.569$ & \textbf{Yes} \\
Temporal Demand  & Ask vs Chat  & 46 & 427.0 & 0.2111  & $-0.210$ & No \\
Frustration      & Form vs Ask  & 52 & 403.5 & 0.0091  & $-0.414$ & \textbf{Yes} \\
Frustration      & Form vs Chat & 57 & 519.5 & 0.0144  & $-0.371$ & \textbf{Yes} \\
Frustration      & Ask vs Chat  & 52 & 629.0 & 0.5838  & $+0.087$ & No \\
AI Reuse         & Form vs Ask  & 49 & 402.5 & 0.0339  & $+0.343$ & No \\
AI Reuse         & Form vs Chat & 43 & 216.0 & 0.0015  & $+0.543$ & \textbf{Yes} \\
AI Reuse         & Ask vs Chat  & 45 & 454.5 & 0.4684  & $+0.122$ & No \\
\bottomrule
\end{tabular}
\caption{Pairwise Wilcoxon signed-rank tests for the survey measures with at least one significant pairwise effect. Sign convention: positive $r$ means the first-named mode scores higher on the measure. Bonferroni-corrected $\alpha = 0.0167$.}
\label{tab:survey-wilcoxon}
\end{table}

\paragraph{Mode preference ranking ($n = 69$).} The chi-square goodness-of-fit on first-rank votes against a uniform null gives $\chi^2(2, N{=}69) = 14.17$, $p < 0.001$, Cram\'er's $V = 0.320$. The Friedman test on the full rank scores (1--3) gives $\chi^2(2) = 5.07$, $p = 0.0792$, $W = 0.037$. Per-mode rank-1 vote counts: Form~37, Ask~12, Chat~20; mean ranks (lower~=~better): Form~1.78, Chat~2.07, Ask~2.14.

\section{NLP and Judge Inferential Tables}
\label{app:stats-nlp}

NLP and judge tests follow the same family: Friedman omnibus on the strict matched-triple $n=26$, Wilcoxon signed-rank pairwise post-hoc with Bonferroni-corrected $\alpha = 0.0167$, plus a Kruskal--Wallis sensitivity check on the full unbalanced $n = 113$. Shapiro--Wilk testing identified non-normal distributions in BLEU, BERTScore Recall (Ask, Chat), and SBERT (Form, Ask); non-parametric tests are used throughout.

\begin{table}[htbp]
\centering
\small
\begin{tabular}{lrrll}
\toprule
\textbf{Metric} & \textbf{H} & \textbf{p} & \textbf{Sig.} & $\boldsymbol{\varepsilon^2}$ \\
\midrule
BLEU                  & 4.48  & 0.107   & No  & --- \\
ROUGE-1               & 5.57  & 0.062   & No  & --- \\
ROUGE-2               & 9.56  & 0.008   & Yes & 0.069 (mod.) \\
ROUGE-L               & 7.79  & 0.020   & Yes & --- \\
METEOR                & 9.80  & 0.007   & Yes & 0.071 (mod.) \\
BERTScore Precision   & 9.11  & 0.011   & Yes & --- \\
BERTScore Recall      & 32.38 & $<$0.001 & Yes & 0.276 (large) \\
BERTScore F1          & 9.75  & 0.008   & Yes & 0.071 (mod.) \\
SBERT                 & 17.43 & $<$0.001 & Yes & 0.140 (large) \\
\bottomrule
\end{tabular}
\caption{Kruskal--Wallis sensitivity check on the full unbalanced sample ($n = 113$).}
\label{tab:nlp-kruskal}
\end{table}

\begin{table}[htbp]
\centering
\small
\begin{tabular}{lrrll}
\toprule
\textbf{Metric} & $\boldsymbol{\chi^2}$ & \textbf{p} & \textbf{Sig.} & $\boldsymbol{W}$ \\
\midrule
BLEU                  & 6.23  & 0.044   & Yes & 0.083 (weak) \\
ROUGE-1               & 10.69 & 0.005   & Yes & $\sim$0.18 (mod.) \\
ROUGE-2               & 13.46 & 0.001   & Yes & $\sim$0.22 (mod.) \\
ROUGE-L               & 11.31 & 0.004   & Yes & $\sim$0.19 (mod.) \\
METEOR                & 9.92  & 0.007   & Yes & $\sim$0.17 (mod.) \\
BERTScore Precision   & 18.54 & $<$0.001 & Yes & $\sim$0.30 (strong) \\
BERTScore Recall      & 30.77 & $<$0.001 & Yes & 0.410 (strong) \\
BERTScore F1          & 13.00 & 0.002   & Yes & $\sim$0.21 (mod.) \\
SBERT                 & 16.69 & $<$0.001 & Yes & $\sim$0.27 (mod.) \\
\bottomrule
\end{tabular}
\caption{Friedman omnibus on the strict within-subjects sample ($n = 26$).}
\label{tab:nlp-friedman}
\end{table}

\begin{table}[htbp]
\centering
\small
\begin{tabular}{llrrrl}
\toprule
\textbf{Metric} & \textbf{Pair} & \textbf{W} & \textbf{p} & \textbf{r} & \textbf{Sig.\ (Bonf.)} \\
\midrule
BLEU                & Form vs Ask  & 108 & 0.089   & $+0.385$ & No \\
BLEU                & Form vs Chat & 133 & 0.291   & $-0.242$ & No \\
BLEU                & Ask vs Chat  & 74  & 0.009   & $-0.578$ & \textbf{Yes} \\
ROUGE-1             & Form vs Ask  & 77  & 0.011   & $+0.561$ & \textbf{Yes} \\
ROUGE-1             & Form vs Chat & 147 & 0.483   & $-0.162$ & No \\
ROUGE-1             & Ask vs Chat  & 63  & 0.003   & $-0.641$ & \textbf{Yes} \\
ROUGE-2             & Form vs Ask  & 30  & $<$0.001 & $+0.829$ & \textbf{Yes} \\
ROUGE-2             & Form vs Chat & 136 & 0.328   & $+0.225$ & No \\
ROUGE-2             & Ask vs Chat  & 71  & 0.007   & $-0.595$ & \textbf{Yes} \\
ROUGE-L             & Form vs Ask  & 44  & $<$0.001 & $+0.749$ & \textbf{Yes} \\
ROUGE-L             & Form vs Chat & 161 & 0.727   & $+0.083$ & No \\
ROUGE-L             & Ask vs Chat  & 66  & 0.004   & $-0.624$ & \textbf{Yes} \\
METEOR              & Form vs Ask  & 33  & $<$0.001 & $+0.812$ & \textbf{Yes} \\
METEOR              & Form vs Chat & 137 & 0.340   & $+0.219$ & No \\
METEOR              & Ask vs Chat  & 84  & 0.019   & $-0.521$ & No \\
BERTScore Precision & Form vs Ask  & 63  & 0.003   & $+0.641$ & \textbf{Yes} \\
BERTScore Precision & Form vs Chat & 106 & 0.080   & $-0.396$ & No \\
BERTScore Precision & Ask vs Chat  & 40  & $<$0.001 & $-0.772$ & \textbf{Yes} \\
BERTScore Recall    & Form vs Ask  & 0   & $<$0.001 & $+1.000$ & \textbf{Yes} \\
BERTScore Recall    & Form vs Chat & 55  & 0.001   & $+0.687$ & \textbf{Yes} \\
BERTScore Recall    & Ask vs Chat  & 99  & 0.053   & $-0.436$ & No \\
BERTScore F1        & Form vs Ask  & 24  & $<$0.001 & $+0.863$ & \textbf{Yes} \\
BERTScore F1        & Form vs Chat & 172 & 0.940   & $+0.020$ & No \\
BERTScore F1        & Ask vs Chat  & 59  & 0.002   & $-0.664$ & \textbf{Yes} \\
SBERT               & Form vs Ask  & 110 & 0.099   & $+0.373$ & No \\
SBERT               & Form vs Chat & 56  & 0.002   & $-0.681$ & \textbf{Yes} \\
SBERT               & Ask vs Chat  & 26  & $<$0.001 & $-0.852$ & \textbf{Yes} \\
\bottomrule
\end{tabular}
\caption{Pairwise Wilcoxon signed-rank tests on the strict within-subjects sample ($n = 26$). Sign convention: positive $r$ means the first-named mode scores higher on the metric. Bonferroni-corrected $\alpha = 0.0167$.}
\label{tab:nlp-wilcoxon}
\end{table}

\paragraph{Confidence--quality correlations.} Spearman $\rho$ between item \texttt{5\_ai} (AI legal-basis confidence) and document quality, computed within each mode on the joined survey~$\times$~NLP sample. Form ($n = 34$): $\rho = -0.22$ ($p = 0.22$) vs.\ BERTScore F1; $\rho = +0.07$ ($p = 0.70$) vs.\ Judge Overall. Ask ($n = 34$): $\rho = -0.23$ ($p = 0.19$) vs.\ BERTScore F1; $\rho = -0.21$ ($p = 0.24$) vs.\ Judge Overall. Chat ($n = 38$): $\rho = +0.28$ ($p = 0.09$) vs.\ BERTScore F1; $\rho = +0.41$ ($p = 0.012$) vs.\ Judge Overall. Pooled ($n = 106$): $\rho = -0.02$ ($p = 0.88$) vs.\ BERTScore F1; $\rho = +0.14$ ($p = 0.14$) vs.\ Judge Overall.
